\documentclass[UTF8,a4paper,12pt]{ctexart}

\usepackage{amsmath}
\usepackage{amssymb}
\usepackage{booktabs}
\usepackage{multirow}
\usepackage{array}
\usepackage{tabularx}
\usepackage[margin=2cm]{geometry}
\usepackage{graphicx}
\usepackage{float}
\usepackage{xcolor}
\usepackage{listings}
\usepackage{caption}
\usepackage{subcaption}
\usepackage{longtable}
\usepackage{hyperref}
\usepackage{enumitem}
\usepackage{makecell}
\usepackage{tikz}
\usetikzlibrary{arrows.meta,positioning}

\hypersetup{
  pdfauthor={},
  pdfsubject={Technical report},
  colorlinks=true,
  linkcolor=black,
  urlcolor=blue,
  citecolor=black
}

\captionsetup{font=small,labelfont=bf}
\setlist[itemize]{leftmargin=2em}
\setlist[enumerate]{leftmargin=2em}
\emergencystretch=2em

\definecolor{codebg}{RGB}{248,248,248}
\definecolor{codekw}{RGB}{22,72,150}
\definecolor{codecm}{RGB}{0,105,95}
\definecolor{codestr}{RGB}{160,90,0}

\lstdefinestyle{pythonstyle}{
  language=Python,
  basicstyle=\ttfamily\small,
  keywordstyle=\color{codekw}\bfseries,
  commentstyle=\color{codecm},
  stringstyle=\color{codestr},
  backgroundcolor=\color{codebg},
  frame=single,
  rulecolor=\color{black!35},
  breaklines=true,
  columns=fullflexible,
  keepspaces=true,
  showstringspaces=false,
  numbers=left,
  numberstyle=\tiny\color{black!45},
  xleftmargin=2em,
  framexleftmargin=1.5em
}
\lstset{style=pythonstyle}

\newcommand{\code}[1]{\texttt{#1}}
\newcommand{\sourceurl}[1]{\footnote{\url{#1}}}
\newcommand{\figplaceholder}[1]{%
  \begin{center}
    \fbox{\parbox{0.78\textwidth}{\centering #1}}
  \end{center}
}
\newcommand{\safeincludegraphics}[2][]{%
  \IfFileExists{#2}{%
    \includegraphics[#1]{#2}%
  }{%
    \figplaceholder{\rule{0pt}{1.2em}}%
  }%
}

\newcommand{\makecover}{
  \begin{titlepage}
    \thispagestyle{empty}
    \begin{center}
      \vspace*{3cm}
      {\LARGE\bfseries 人脸表情识别系统技术报告\par}
      \vspace*{1cm}
      {\Large 基于 PyQt5 与深度学习\par}
      \vspace*{2cm}
      {\large \par}
    \end{center}
  \end{titlepage}
}

\begin{document}

\makecover
\tableofcontents
\newpage

\section{摘要}

本项目面向人脸表情识别任务，构建了一个从模型训练、外部泛化测试、迁移学习实验到桌面端交互展示的完整系统。项目基于 PyTorch 实现了 SimpleCNN、ImprovedCNN、MiniResNet 和 RGBMiniResNet 等卷积神经网络模型，使用 FER2013、CK+、RAF-DB、AffectNet、RAF-DB Compound 和 RAF-ML 等数据集完成了基本七类表情识别、外部泛化测试、RGB 彩色模型训练、AffectNet 预训练迁移、复合表情识别和多标签混合表情识别。

在系统实现方面，项目搭建了基于 PyQt5 的图形界面，支持图片上传识别和摄像头实时识别。GUI 支持基本表情识别、复合表情识别和混合表情识别三类任务，并可动态切换 softmax 单标签输出和 sigmoid 多标签输出。摄像头部分支持 YuNet、MediaPipe 和中心区域三种人脸定位方式，能够在实际演示中显示预测结果、置信度和各类别概率。

实验结果表明，RAF-DB RGBMiniResNet 在 RAF-DB test 上取得 Accuracy 0.8380、Macro F1 0.7591，明显优于 FER2013 灰度模型在 RAF-DB 外部测试上的表现。AffectNet 预训练 + RAF-DB 微调模型取得 Accuracy 0.8181、Macro F1 0.7302，说明迁移学习具有一定价值，但仍受到数据集分布差异影响。RAF-DB Compound 复合表情识别 Macro F1 为 0.4492，体现出复合表情类别样本少、语义接近和类别边界模糊的难度。RAF-ML 多标签混合表情识别 Macro F1 为 0.7352，说明多标签形式更适合表达真实场景中多个情绪成分同时存在的情况。项目最后通过消融实验分析了输入形式、训练策略和任务形式对模型性能的影响。

项目最终形成了包含数据准备、模型训练、模型评估、消融分析和 GUI 展示的完整实验闭环，为课程 Final Project 提供了从算法实现到系统演示的综合实践案例。

\section{实验简介与动机}

人脸表情识别是情感计算和人机交互中的重要研究方向。相比普通图像分类任务，表情识别不仅需要关注整体面部结构，还需要识别嘴角、眉眼、鼻翼、面部肌肉纹理等局部变化。真实应用中，人脸表情还会受到光照、姿态、遮挡、图像清晰度和摄像头质量的影响，因此模型的泛化能力和实时可用性同样重要。

本项目选择人脸表情识别作为人工智能课程 Final Project，主要动机包括以下三点。第一，该任务与深度学习图像分类高度相关，适合从基础 CNN 逐步扩展到残差网络和 RGB 模型。第二，表情识别具有明确的交互展示价值，可以通过 PyQt5 GUI 和摄像头形成可演示系统。第三，真实表情并不总是单一的基本情绪，因此项目进一步扩展到 RAF-DB Compound 复合表情识别和 RAF-ML 多标签混合表情识别，用于分析不同任务定义对模型效果的影响。

本项目为单人项目，由本人独立完成实验设计、数据集选择、模型训练测试、结果分析和 GUI 调试。部分实现思路和代码框架在 AI 辅助下完成，但整体实验路线、数据处理、结果解释和系统调试均由本人完成。

\section{项目环境与代码结构}

\subsection{实验环境}

本项目的训练、评估和 GUI 演示均在统一的软件环境下完成，主要依赖包括 PyTorch、PyQt5、OpenCV、pandas、scikit-learn 和 matplotlib 等。实验环境配置见表~\ref{tab:experiment_env}。

\begin{table}[H]
  \centering
  \caption{项目实验环境}
  \label{tab:experiment_env}
  \begin{tabularx}{\textwidth}{lX}
    \toprule
    项目 & 内容 \\
    \midrule
    操作系统 & Windows \\
    Python 环境 & \code{face\_emotion\_env} \\
    深度学习框架 & PyTorch \\
    GUI 框架 & PyQt5 \\
    图像处理 & OpenCV、Pillow \\
    数据处理 & pandas、numpy \\
    评估与可视化 & scikit-learn、matplotlib \\
    人脸定位 & OpenCV YuNet、MediaPipe Tasks API、中心区域识别 \\
    运行设备 & 支持 CUDA 时使用 GPU，否则使用 CPU \\
    \bottomrule
  \end{tabularx}
\end{table}

\subsection{代码结构}

项目代码按照模型定义、数据处理、训练评估、GUI 展示和结果输出进行组织。主要文件与目录说明见表~\ref{tab:project_structure}。

\small
\begin{longtable}{>{\raggedright\arraybackslash}p{0.43\textwidth}>{\raggedright\arraybackslash}p{0.49\textwidth}}
  \caption{项目主要文件与目录说明}\label{tab:project_structure}\\
  \toprule
  路径或文件 & 作用 \\
  \midrule
  \endfirsthead
  \toprule
  路径或文件 & 作用 \\
  \midrule
  \endhead
  \code{models/} & 存放 SimpleCNN、ImprovedCNN、MiniResNet、RGBMiniResNet 等模型定义。\\
  \code{models/detectors/} & 存放 YuNet ONNX 和 MediaPipe tflite 人脸检测模型。\\
  \code{utils/} & 存放数据集读取、指标计算、预测预处理和可视化等工具函数。\\
  \code{train.py} & FER2013 基础灰度模型训练脚本。\\
  \code{evaluate.py} & FER2013 基础模型评估脚本。\\
  \code{ckplus\_eval.py} & CK+ 外部泛化测试脚本。\\
  \code{rafdb\_eval.py} & FER2013 灰度模型在 RAF-DB 上的外部测试脚本。\\
  \code{train\_rafdb\_rgb.py} & RAF-DB RGBMiniResNet 训练脚本。\\
  \code{evaluate\_rafdb\_rgb.py} & RAF-DB RGB 模型评估脚本。\\
  \code{train\_affectnet\_rgb.py} & AffectNet RGB 模型训练脚本。\\
  \code{evaluate\_affectnet\_rgb.py} & AffectNet RGB 模型测试脚本。\\
  \code{finetune\_rafdb\_from\_affectnet.py} & AffectNet 预训练后在 RAF-DB 微调的脚本。\\
  \code{train\_compound\_rgb.py} & RAF-DB Compound 复合表情训练脚本。\\
  \code{evaluate\_compound\_rgb.py} & RAF-DB Compound 复合表情评估脚本。\\
  \code{train\_rafml\_multilabel.py} & RAF-ML 多标签混合表情训练脚本。\\
  \code{evaluate\_rafml\_multilabel.py} & RAF-ML 多标签混合表情评估脚本。\\
  \code{analyze\_ablation\_results.py} & 消融实验结果汇总与图表生成脚本。\\
  \code{gui.py} & PyQt5 图形界面系统。\\
  \code{checkpoints/} & 保存模型最优权重。\\
  \code{logs/} & 保存训练日志、loss 曲线、accuracy 或 F1 曲线。\\
  \code{outputs/} & 保存分类报告、混淆矩阵、类别统计、样例和消融分析结果。\\
  \bottomrule
\end{longtable}
\normalsize

\section{数据集介绍与预处理}

本项目使用的数据集来源较多，既包括 Kaggle 上整理后的 ImageFolder 或 YOLO Format 数据，也包括 RAF 官方页面介绍的 Basic、Compound 和 Multi-label 数据。为避免不同数据集标签顺序不一致导致实验结果不可比，项目将基本表情任务统一映射到 FER2013 顺序：Angry、Disgust、Fear、Happy、Sad、Surprise、Neutral。各数据集与任务对应关系见表~\ref{tab:dataset_task_map}。

\begin{table}[H]
  \centering
  \caption{数据集与任务对应关系}
  \label{tab:dataset_task_map}
  \scriptsize
  \setlength{\tabcolsep}{3pt}
  \begin{tabularx}{\textwidth}{p{0.13\textwidth}p{0.16\textwidth}p{0.13\textwidth}p{0.09\textwidth}p{0.17\textwidth}X}
    \toprule
    数据集 & 来源 & 输入形式 & 类别数 & 任务类型 & 项目用途 \\
    \midrule
    FER2013 & Kaggle ImageFolder & 灰度 48$\times$48 & 7 & softmax 单标签 & 基础灰度模型训练。\\
    CK+ & Kaggle ckdataset & 灰度 48$\times$48 & 7 & softmax 单标签 & 外部泛化测试，删除 Contempt。\\
    RAF-DB Basic & RAF 官方 / Kaggle & RGB 96$\times$96 & 7 & softmax 单标签 & RGB 主模型训练与测试。\\
    AffectNet YOLO Format & Kaggle YOLO Format & RGB 96$\times$96 & 7 & softmax 单标签 & 真实场景训练与预训练。\\
    RAF-DB Compound & RAF 官方 & RGB 96$\times$96 & 11 & softmax 单标签 & 复合表情扩展实验。\\
    RAF-ML & RAF 官方 & RGB 96$\times$96 & 6 & sigmoid 多标签 & 混合表情扩展实验。\\
    \bottomrule
  \end{tabularx}
\end{table}

FER2013 和 CK+ 都以 48$\times$48 灰度图像形式用于基础模型训练或外部泛化测试，二者的数据说明图见图~\ref{fig:gray_dataset_intro}。

\begin{figure}[H]
  \centering
  \begin{subfigure}{0.48\textwidth}
    \centering
    \figplaceholder{\rule{0pt}{1.2em}}
    \caption{FER2013 数据说明图}
  \end{subfigure}
  \hfill
  \begin{subfigure}{0.48\textwidth}
    \centering
    \figplaceholder{\rule{0pt}{1.2em}}
    \caption{CK+ 数据说明图}
  \end{subfigure}
  \caption{灰度基础数据集说明图}
  \label{fig:gray_dataset_intro}
\end{figure}

图~\ref{fig:gray_dataset_intro} 展示了灰度基础数据的典型来源。FER2013 更适合作为训练基线，CK+ 则用于检验灰度模型在较清晰实验室表情数据上的外部泛化能力。

RAF-DB、AffectNet 和 RAF-ML 面向更复杂的真实场景彩色或混合表情任务，其数据说明图见图~\ref{fig:rgb_dataset_intro}。

\begin{figure}[H]
  \centering
  \begin{subfigure}{0.32\textwidth}
    \centering
    \figplaceholder{\rule{0pt}{1.2em}}
    \caption{RAF-DB}
  \end{subfigure}
  \hfill
  \begin{subfigure}{0.32\textwidth}
    \centering
    \figplaceholder{\rule{0pt}{1.2em}}
    \caption{AffectNet YOLO}
  \end{subfigure}
  \hfill
  \begin{subfigure}{0.32\textwidth}
    \centering
    \figplaceholder{\rule{0pt}{1.2em}}
    \caption{RAF-ML}
  \end{subfigure}
  \caption{真实场景彩色与混合表情数据集说明图}
  \label{fig:rgb_dataset_intro}
\end{figure}

从图~\ref{fig:rgb_dataset_intro} 可以看出，真实场景数据在图像风格、表情强度和任务定义上比灰度基础数据更加复杂，这也是本项目进一步引入 RGB 模型、复合表情任务和多标签混合表情任务的原因。

\subsection{FER2013}

FER2013 是人脸表情识别领域常用的经典基线数据集，图像为低分辨率灰度人脸，标签定义清晰，适合用于基础表情分类模型训练和算法对比。由于 FER2013 图像分辨率较低且不包含颜色信息，它也可以用于观察灰度模型在真实彩色场景中的跨数据集泛化限制。

本项目使用 Kaggle 上整理为 ImageFolder 结构的 FER2013 版本\sourceurl{https://www.kaggle.com/datasets/astraszab/facial-expression-dataset-image-folders-fer2013}，实际路径为 \code{data/fer2013}。该数据集被划分为 train、test 和 validation，标签为 0 到 6，分别对应 Angry、Disgust、Fear、Happy、Sad、Surprise、Neutral。项目中 FER2013 用于训练 SimpleCNN、ImprovedCNN 和灰度 MiniResNet，输入图像统一处理为 48$\times$48 灰度图，并进行归一化和必要的数据增强。

\subsection{CK+}

CK+ 是实验室环境下采集的人脸表情序列数据集，表情通常较清晰、较夸张，常用于表情识别算法验证。与真实网络图像相比，CK+ 的采集条件更受控，因此适合作为外部测试集观察模型在较清晰表情样本上的泛化表现。

本项目使用 Kaggle 上整理的 ckdataset\sourceurl{https://www.kaggle.com/datasets/davilsena/ckdataset}，实际文件为 \code{data/ck\_dataset/ckextended.csv}。该文件包含 \code{emotion}、\code{pixels}、\code{Usage} 三列，其中每行 \code{pixels} 包含 2304 个空格分隔像素值，对应 48$\times$48 灰度人脸图像。原始标签包含 8 类：Anger、Disgust、Fear、Happiness、Sadness、Surprise、Neutral 和 Contempt。项目中删除 emotion=7 的 Contempt 类，仅保留七类基本表情，用于测试 FER2013 灰度 MiniResNet 的外部泛化能力。

\subsection{RAF-DB Basic}

RAF-DB 是 Real-world Affective Faces Database，官方页面说明其包含约 30K 张从互联网收集的真实场景人脸图像，并由约 40 名标注者进行众包标注\sourceurl{http://www.whdeng.cn/RAF/model1.html}。与 FER2013 和 CK+ 相比，RAF-DB 具有更明显的真实应用特征，包括姿态、光照、遮挡、年龄、性别、种族和图像后处理差异，因此更适合作为真实场景 RGB 表情识别主实验数据。

本项目实际使用 Kaggle 上整理的 RAF-DB Basic ImageFolder 版本\sourceurl{https://www.kaggle.com/datasets/shuvoalok/raf-db-dataset}，路径为 \code{data/RAF-DB/DATASET}，其中 \code{train/1..7} 和 \code{test/1..7} 分别存放七类基本表情图像。basic 部分共 15339 张图像，其中 train 12271 张、test 3068 张，用于 RGBMiniResNet 主模型训练、测试，以及与灰度外部测试进行输入形式消融对比。

RAF-DB 的 basic 标签顺序为 Surprise、Fear、Disgust、Happy、Sad、Angry、Neutral，而本项目统一使用 FER2013 顺序。因此代码中使用 \code{RAF\_TO\_FER} 完成标签映射，如 Listing~\ref{lst:label_mapping} 所示。

\begin{lstlisting}[language=Python, caption={RAF-DB 与 AffectNet 标签映射核心代码}, label={lst:label_mapping}]
RAF_TO_FER = {
    1: 5,  # Surprise -> Surprise
    2: 2,  # Fear -> Fear
    3: 1,  # Disgust -> Disgust
    4: 3,  # Happy -> Happy
    5: 4,  # Sad -> Sad
    6: 0,  # Angry -> Angry
    7: 6,  # Neutral -> Neutral
}

AFFECTNET_TO_FER = {
    0: 0,  # Anger -> Angry
    2: 1,  # Disgust -> Disgust
    3: 2,  # Fear -> Fear
    4: 3,  # Happy -> Happy
    6: 4,  # Sad -> Sad
    7: 5,  # Surprise -> Surprise
    5: 6,  # Neutral -> Neutral
}
\end{lstlisting}

Listing~\ref{lst:label_mapping} 将 RAF-DB 和 AffectNet 的原始标签统一到 FER2013 七类顺序，为后续跨数据集评估、RGB 训练和 GUI 统一显示概率条提供了基础。

\subsection{AffectNet YOLO Format}

AffectNet 是大规模真实场景表情数据集，图像来源复杂，表情变化更丰富，适合用于真实场景表情特征学习和预训练。相比 RAF-DB Basic，AffectNet 的图像风格更分散，标注噪声和类别分布差异也更明显，因此它既能提供更丰富的训练样本，也可能带来跨数据集迁移时的分布差异。

本项目没有使用原始 AffectNet 官网数据，而是使用 Kaggle 上整理好的 AffectNet YOLO Format 版本\sourceurl{https://www.kaggle.com/datasets/fatihkgg/affectnet-yolo-format}。项目实际路径为 \code{data/AffectNet/YOLO\_format}，包含 \code{train/images}、\code{train/labels}、\code{valid/images}、\code{valid/labels}、\code{test/images}、\code{test/labels} 和 \code{data.yaml}。图像按 96$\times$96 RGB 输入处理，原始 8 类为 Anger、Contempt、Disgust、Fear、Happy、Neutral、Sad、Surprise。项目删除 Contempt 类，并映射为统一七类基本表情，用于 AffectNet RGB Full 模型训练和 AffectNet 预训练 + RAF-DB 微调实验。

\subsection{RAF-DB Compound}

RAF-DB Compound 是 RAF 系列中的复合表情子集，使用固定的 11 类复合表情表示复杂情绪，例如 Happily Surprised、Sadly Fearful 和 Angrily Disgusted 等。与基本七类表情相比，Compound 类别语义更接近、边界更模糊，并且样本数量更少，因此识别难度明显更高。

本实验使用 \code{data/compound} 下的 aligned 图像和标签文件，共 3954 张图像，其中 train 3162 张、test 792 张。该任务将复合情绪视为 11 个互斥类别，使用 RGB 96$\times$96 输入和 softmax 单标签分类，用于扩展分析固定复合类别任务的难度。类别如表 \ref{tab:compound_classes} 所示。

\begin{table}[H]
  \centering
  \caption{RAF-DB Compound 11 类复合表情}
  \label{tab:compound_classes}
  \small
  \begin{tabular}{cll}
    \toprule
    编号 & 英文类别 & 中文含义 \\
    \midrule
    1 & Happily Surprised & 开心-惊讶 \\
    2 & Happily Disgusted & 开心-厌恶 \\
    3 & Sadly Fearful & 悲伤-恐惧 \\
    4 & Sadly Angry & 悲伤-愤怒 \\
    5 & Sadly Surprised & 悲伤-惊讶 \\
    6 & Sadly Disgusted & 悲伤-厌恶 \\
    7 & Fearfully Angry & 恐惧-愤怒 \\
    8 & Fearfully Surprised & 恐惧-惊讶 \\
    9 & Angrily Surprised & 愤怒-惊讶 \\
    10 & Angrily Disgusted & 愤怒-厌恶 \\
    11 & Disgustedly Surprised & 厌恶-惊讶 \\
    \bottomrule
  \end{tabular}
\end{table}

表~\ref{tab:compound_classes} 中的 11 个类别均为固定复合表情类别，模型需要在语义相近的组合情绪之间进行区分，因此该任务比基本七类表情识别更困难。

\subsection{RAF-ML}

RAF-ML 是 Real-world Affective Faces Multi Label 数据集，官方页面说明其包含 4908 张真实场景混合表情图像，每张图像具有 6 维表达分布向量和多标签情绪标注\sourceurl{http://www.whdeng.cn/RAF/model2.html}。与 Compound 将混合表情压缩为固定互斥类别不同，RAF-ML 用 Surprise、Fear、Disgust、Happiness、Sadness、Anger 六维标签表示多个情绪成分是否同时存在，不包含 Neutral。

本实验使用 \code{data/RAF-ML}，读取 RAF-ML 图像及多标签标注文件，输入为 RGB 96$\times$96。模型输出 6 维 sigmoid 多标签概率，训练损失使用 BCEWithLogitsLoss，评估时主要使用 Micro F1 和 Macro F1，而不使用普通单标签 Accuracy。

\subsection{统一预处理}

不同任务的输入预处理有所区别。灰度模型使用 48$\times$48 单通道输入，RGB 模型使用 96$\times$96 三通道输入。项目对不同数据集进行了标签映射、Contempt 类删除、图像 resize、归一化和数据增强。外部测试时，项目保持测试数据处理与模型训练输入格式一致，避免输入尺寸和标签顺序不匹配。

\subsection{数据集规模与类别分布分析}

表情识别数据集普遍存在类别不平衡问题。多数真实数据集中 Happy、Neutral 等类别样本更容易收集，而 Fear、Disgust 或部分复合表情类别样本较少。类别不平衡会使 Accuracy 更容易受到多数类影响，因此本项目在报告 Accuracy 的同时，也使用 Macro F1、Micro F1、每类 Precision/Recall/F1 和混淆矩阵进行分析。主要数据集在测试或评估阶段的类别规模见表~\ref{tab:dataset_distribution}。

\begin{table}[H]
  \centering
  \caption{主要数据集类别规模与分布特点}
  \label{tab:dataset_distribution}
  \scriptsize
  \setlength{\tabcolsep}{3pt}
  \renewcommand{\arraystretch}{1.15}
  \begin{tabularx}{\textwidth}{p{0.16\textwidth}p{0.20\textwidth}p{0.27\textwidth}X}
    \toprule
    数据集 & 类别或标签规模 & 已有统计示例 & 分布特点与影响 \\
    \midrule
    FER2013 & 7 类基本表情 & 测试集：Happy 879，Disgust 55，Neutral 626，Sad 594，Fear 528，Angry 491，Surprise 416 & Happy 类样本相对较多，Disgust 类明显较少，小类 Recall 和 Macro F1 容易受到影响。\\
    RAF-DB Basic & 7 类基本表情 & 测试集：Happy 1185，Neutral 680，Sad 478，Surprise 329，Angry 162，Disgust 160，Fear 74 & Happy、Neutral 等类别数量较多，Fear、Disgust 较少，模型容易更偏向多数类。\\
    RAF-DB Compound & 11 类复合表情 & 测试集支持数范围为 18--174，例如 Angrily Disgusted 174，Sadly Surprised 18 & 总量小于 Basic，类别更细且分布更不均衡，复合类别的 Macro F1 更容易下降。\\
    RAF-ML & 6 维多标签 & 测试集正样本：Surprise 408，Fear 232，Disgust 539，Happiness 258，Sadness 343，Anger 384 & 每个标签正样本数不同，Micro F1 与 Macro F1 的差异能够反映标签分布和标签难度差异。\\
    \bottomrule
  \end{tabularx}
\end{table}

从表~\ref{tab:dataset_distribution} 可以看出，不同数据集的不平衡形式并不相同。FER2013 和 RAF-DB Basic 的不平衡主要体现在基本表情类别数量差异；RAF-DB Compound 的难点则叠加了小样本、细粒度和类别语义接近；RAF-ML 虽然采用多标签形式，但各标签正样本数不同，因此需要同时观察 Micro F1 和 Macro F1。后续实验分析中，混淆矩阵和每类指标用于进一步解释这些分布差异带来的影响。

\section{模型设计}

\subsection{模型设计总体思路}

本项目的模型设计不是围绕单一网络展开，而是从基础灰度 CNN 逐步扩展到 RGB 残差模型，并进一步适配复合表情和多标签混合表情任务。整体路线可以概括为三层：第一层使用 FER2013 灰度图像训练 SimpleCNN、ImprovedCNN 和 MiniResNet，建立基本七类表情识别基线；第二层将输入从 48$\times$48 灰度图扩展到 96$\times$96 RGB 图像，在 RAF-DB 和 AffectNet 上训练 RGBMiniResNet；第三层复用 RGBMiniResNet 主体结构，通过改变输出维度和损失函数，支持 RAF-DB Compound 11 类复合表情识别和 RAF-ML 6 维多标签混合表情识别。

这样的设计使模型复杂度随任务逐步增加：SimpleCNN 便于快速验证训练评估流程，ImprovedCNN 用于观察 BatchNorm、Dropout 和数据增强的常规改进效果，MiniResNet 引入残差连接提升灰度模型表达能力，RGBMiniResNet 则面向真实彩色图像和 GUI 摄像头演示。三类任务的输出形式不同：基本表情识别和复合表情识别是互斥类别分类，使用 softmax；RAF-ML 混合表情识别允许多个情绪成分同时存在，使用 sigmoid 多标签输出。

\subsection{SimpleCNN}

SimpleCNN 是项目中最基础的灰度卷积神经网络，输入为 $1\times48\times48$，输出为 7 类基本表情 logits，训练时配合 CrossEntropyLoss 和 softmax 单标签分类。其主要结构为两组卷积、ReLU 和最大池化，随后通过 Flatten 和全连接层得到类别输出。该模型没有使用 BatchNorm、Dropout 或残差连接，因此结构最简单，训练速度也最快。

SimpleCNN 的作用是建立 FER2013 灰度表情识别基线，用于确认数据读取、训练循环、日志保存、分类报告和混淆矩阵生成流程是否正确。它的优点是参数量较少、实现清晰、便于调试；缺点是特征表达能力有限，对姿态、光照和跨数据集风格变化的鲁棒性较弱。因此在后续实验中，SimpleCNN 主要作为对照模型，而不是最终推荐模型。

\subsection{ImprovedCNN}

ImprovedCNN 在 SimpleCNN 的基础上加深网络结构，输入仍为 $1\times48\times48$ 灰度图，输出仍为 7 类 softmax 单标签分类。与 SimpleCNN 相比，它使用 32、64、128 通道的多层卷积结构，并在卷积后加入 BatchNorm 和 ReLU，以稳定训练过程、缓解梯度分布变化。模型中还加入 Dropout2d 和全连接层 Dropout，用于抑制过拟合。

该模型的设计目的，是验证普通 CNN 在加入常见正则化和归一化技巧后，是否能够比基础 CNN 获得更好的 FER2013 表现。BatchNorm 可以让训练过程更平稳，Dropout 可以降低模型对训练集局部模式的依赖，数据增强则提高模型对轻微平移、翻转和亮度变化的鲁棒性。ImprovedCNN 的优点是仍然容易训练，但表达能力强于 SimpleCNN；缺点是缺少显式残差连接，进一步加深时训练稳定性和泛化能力仍然有限。

\subsection{MiniResNet}

MiniResNet 是面向 48$\times$48 灰度表情图像的轻量残差网络，输入为单通道灰度图，输出为 7 类基本表情 logits。模型先经过一个 stem 卷积层，再经过多个 ResidualBlock，通道数从 32 逐步提升到 64 和 128，最后使用 AdaptiveAvgPool2d 将空间特征聚合为 $1\times1$，再接线性分类层。

残差连接的核心思想是让网络学习残差映射而不是完整映射，即将主分支卷积输出与 shortcut 分支相加。这样可以缓解较深网络训练中的梯度传播困难，使模型在增加深度时仍能保持稳定训练。与 ImprovedCNN 相比，MiniResNet 参数量仍然较小，但通过残差块获得了更好的特征表达能力，因此被作为 FER2013 灰度主实验模型，并用于 CK+ 外部测试和 RAF-DB 灰度外部测试。

MiniResNet 的优点是结构轻量、训练成本低、性能优于普通浅层 CNN，适合课程项目和本地训练环境；缺点是输入仍为灰度低分辨率图像，无法利用真实彩色照片中的颜色和纹理信息，因此跨到 RAF-DB 等真实 RGB 数据集时泛化能力有限。

\subsection{RGBMiniResNet}

由于 RAF-DB、AffectNet、Compound 和 RAF-ML 均使用 RGB 图像，本项目在灰度 MiniResNet 的基础上实现了 RGBMiniResNet。其输入从 1 通道扩展到 3 通道，尺寸从 48$\times$48 扩展到 96$\times$96，并增加到四个残差阶段，通道数依次为 32、64、128 和 192。更大的输入尺寸和 RGB 信息让模型能够利用颜色、纹理和更细粒度的面部局部变化。

RGBMiniResNet 的核心是残差块和自适应池化分类头。前者用于增强特征表达，后者使用 AdaptiveAvgPool2d 降低全连接层参数量，并使分类头只依赖最终通道数而不是固定空间尺寸。分类头中的 Linear 层通过 \code{num\_classes} 动态设置输出维度，因此同一主干结构可以复用于 7 类基本表情、11 类复合表情和 6 维多标签混合表情任务。核心结构如 Listing~\ref{lst:rgb_resnet} 所示。

\begin{lstlisting}[language=Python, caption={RGBMiniResNet 残差块与分类头核心结构}, label={lst:rgb_resnet}]
class RGBResidualBlock(nn.Module):
    def __init__(self, in_channels: int, out_channels: int, stride: int = 1):
        super().__init__()
        self.conv = nn.Sequential(
            nn.Conv2d(in_channels, out_channels, 3, stride=stride, padding=1, bias=False),
            nn.BatchNorm2d(out_channels),
            nn.ReLU(inplace=True),
            nn.Conv2d(out_channels, out_channels, 3, padding=1, bias=False),
            nn.BatchNorm2d(out_channels),
        )
        if stride != 1 or in_channels != out_channels:
            self.shortcut = nn.Sequential(
                nn.Conv2d(in_channels, out_channels, 1, stride=stride, bias=False),
                nn.BatchNorm2d(out_channels),
            )
        else:
            self.shortcut = nn.Identity()
        self.relu = nn.ReLU(inplace=True)

    def forward(self, x: torch.Tensor) -> torch.Tensor:
        return self.relu(self.conv(x) + self.shortcut(x))

self.pool = nn.AdaptiveAvgPool2d((1, 1))
self.classifier = nn.Sequential(
    nn.Dropout(0.25),
    nn.Linear(192, num_classes),
)
\end{lstlisting}

代码中 \code{self.shortcut} 用于在输入输出通道数或步长不一致时进行维度匹配；\code{forward} 函数将主分支卷积结果与 shortcut 相加后再经过 ReLU，形成残差连接。分类头使用 AdaptiveAvgPool2d 将空间特征压缩为 $1\times1$，再通过 Dropout 和 Linear 输出对应类别数。由于 Linear 的输出维度由 \code{num\_classes} 决定，因此同一结构可以复用于 7 类、11 类和 6 维多标签任务。

RGBMiniResNet 的优点是能够处理真实彩色人脸图像，并且通过轻量残差结构在训练成本和表达能力之间取得平衡；缺点是主干网络仍然较轻，与 ResNet18、MobileNet、EfficientNet 或 Vision Transformer 等更强 backbone 相比，复杂场景下仍有提升空间。

\subsection{不同任务下的输出头设计}

基本表情识别和复合表情识别均采用 softmax 单标签分类。前者输出 7 类基本表情，后者输出 11 类复合表情，损失函数使用 CrossEntropyLoss。RAF-ML 混合表情识别采用 sigmoid 多标签输出，输出 6 个独立情绪成分概率，损失函数使用 BCEWithLogitsLoss。

其中 basic\_7 和 compound\_11 都假设一张图像只属于一个互斥类别，因此 softmax 会将所有类别概率归一化为和为 1。RAF-ML 则不同，一张图像可能同时具有 Happiness 和 Surprise 等多个情绪成分，如果使用 softmax 会强制多个标签互相竞争，破坏多标签语义。因此 RAF-ML 使用 sigmoid 对每个情绪维度独立建模，并根据阈值判断哪些情绪成分同时存在。三类任务的输出方式对比见表~\ref{tab:output_types}。

\begin{table}[H]
  \centering
  \caption{三类任务输出方式对比}
  \label{tab:output_types}
  \begin{tabularx}{\textwidth}{lcccX}
    \toprule
    任务 & 输出方式 & 类别数 & 损失函数 & 说明 \\
    \midrule
    基本表情识别 & softmax & 7 & CrossEntropyLoss & 每张图像只属于一个基本表情类别。\\
    复合表情识别 & softmax & 11 & CrossEntropyLoss & 每张图像归为一个固定复合表情类别。\\
    混合表情识别 & sigmoid & 6 & BCEWithLogitsLoss & 每张图像可以同时包含多个情绪成分。\\
    \bottomrule
  \end{tabularx}
\end{table}

表~\ref{tab:output_types} 说明了三类任务在输出方式和损失函数上的根本差异，后续训练和 GUI 预测均围绕该差异自动选择 softmax 或 sigmoid 处理方式。

\subsection{模型版本与权重对应关系}

为了避免总表过宽，报告将模型与权重对应关系拆成灰度基础模型和 RGB/扩展任务模型两张表。灰度基础模型与权重文件的对应关系见表~\ref{tab:gray_model_weight_map}，RGB 与扩展任务模型的权重对应关系见表~\ref{tab:rgb_model_weight_map}。灰度模型均服务于 FER2013 基础实验及外部泛化测试；RGB 模型则覆盖 RAF-DB、AffectNet、Compound 和 RAF-ML，并在 GUI 中作为主要可选模型。

\begin{table}[H]
  \centering
  \caption{灰度基础模型与权重对应关系}
  \label{tab:gray_model_weight_map}
  \small
  \setlength{\tabcolsep}{3pt}
  \begin{tabularx}{\textwidth}{p{0.18\textwidth}p{0.16\textwidth}p{0.13\textwidth}p{0.13\textwidth}p{0.20\textwidth}X}
    \toprule
    模型显示名称 & 模型类 & 输入 & 输出 & 权重文件 & 用途 \\
    \midrule
    SimpleCNN & \code{SimpleCNN} & 灰度 48$\times$48 & 7 类 softmax 单标签 & \path{checkpoints/best_simple_cnn.pth} & FER2013 基线模型，验证基础训练流程。\\
    ImprovedCNN & \code{ImprovedCNN} & 灰度 48$\times$48 & 7 类 softmax 单标签 & \path{checkpoints/best_improved_cnn.pth} & 加入 BatchNorm、Dropout 和数据增强的改进灰度模型。\\
    FER2013 灰度 MiniResNet & \code{MiniResNet} & 灰度 48$\times$48 & 7 类 softmax 单标签 & \path{checkpoints/best_mini_resnet.pth} & FER2013 主实验、CK+ 外部测试和 RAF-DB 灰度外部测试。\\
    \bottomrule
  \end{tabularx}
\end{table}

表~\ref{tab:gray_model_weight_map} 体现了灰度模型从基础 CNN 到残差网络的递进关系，三者均面向七类基本表情识别，但网络复杂度和泛化能力不同。

\begin{table}[H]
  \centering
  \caption{RGB 与扩展任务模型权重对应关系}
  \label{tab:rgb_model_weight_map}
  \scriptsize
  \setlength{\tabcolsep}{3pt}
  \begin{tabularx}{\textwidth}{p{0.22\textwidth}p{0.12\textwidth}p{0.13\textwidth}p{0.14\textwidth}p{0.22\textwidth}X}
    \toprule
    模型显示名称 & 模型类 & 输入 & 输出/损失 & 权重文件 & 用途 \\
    \midrule
    RAF-DB RGBMiniResNet & \code{RGBMiniResNet} & RGB 96$\times$96 & 7 类 softmax 单标签 / CrossEntropyLoss & \path{checkpoints/best_rafdb_rgb_mini_resnet.pth} & RAF-DB Basic 同源训练与测试。\\
    AffectNet RGB Full MiniResNet & \code{RGBMiniResNet} & RGB 96$\times$96 & 7 类 softmax 单标签 / CrossEntropyLoss & \path{checkpoints/best_affectnet_rgb_full_40epoch.pth} & AffectNet 全量训练，GUI 摄像头默认推荐模型。\\
    AffectNet 预训练 + RAF-DB 微调 & \code{RGBMiniResNet} & RGB 96$\times$96 & 7 类 softmax 单标签 / CrossEntropyLoss & \path{checkpoints/best_affectnet_pretrained_rafdb_rgb_mini_resnet.pth} & 迁移学习对比实验。\\
    Compound 11 类 RGBMiniResNet & \code{RGBMiniResNet} & RGB 96$\times$96 & 11 类 softmax 单标签 / CrossEntropyLoss & \path{checkpoints/best_compound_rgb_affectnet_pretrained.pth} & RAF-DB Compound 复合表情识别。\\
    RAF-ML 6 维多标签 RGBMiniResNet & \code{RGBMiniResNet} & RGB 96$\times$96 & 6 维 sigmoid 多标签 / BCEWithLogitsLoss & \path{checkpoints/best_rafml_affectnet_pretrained.pth} & RAF-ML 多标签混合表情识别。\\
    \bottomrule
  \end{tabularx}
\end{table}

表~\ref{tab:rgb_model_weight_map} 展示了 RGBMiniResNet 在不同任务中的复用方式：当输出维度为 7 或 11 时执行 softmax 单标签分类，当输出维度为 6 时执行 sigmoid 多标签识别。

\subsection{评价指标}

本项目使用 Accuracy、Precision、Recall、F1-score、Macro F1 和 Micro F1 等指标评价模型。单标签任务主要关注 Accuracy 和 Macro F1，多标签任务主要关注 Micro F1 和 Macro F1。

\begin{align}
  \mathrm{Accuracy} &= \frac{TP+TN}{TP+TN+FP+FN},\\
  \mathrm{Precision} &= \frac{TP}{TP+FP},\\
  \mathrm{Recall} &= \frac{TP}{TP+FN},\\
  \mathrm{F1} &= \frac{2\cdot \mathrm{Precision}\cdot \mathrm{Recall}}{\mathrm{Precision}+\mathrm{Recall}}.
\end{align}

Macro F1 对每个类别分别计算 F1 后取平均，更能反映类别不平衡场景下的小类表现；Micro F1 汇总所有类别的 TP、FP 和 FN 后计算，更适合多标签任务的整体性能评价。

\section{实验设计与训练流程}

本项目围绕人脸表情识别任务，按照数据准备、统一预处理、模型训练、模型评估、消融对比和 GUI 展示等环节展开。整体实验流程如图~\ref{fig:experiment_pipeline} 所示。

\begin{figure}[H]
  \centering
  \begin{tikzpicture}[scale=0.82, transform shape,
    node distance=1.1cm,
    box/.style={draw, rounded corners, align=center, minimum width=3.1cm, minimum height=0.9cm, fill=blue!5},
    arrow/.style={-{Stealth}, thick}
  ]
    \node[box] (data) {多数据集输入\\FER2013 / RAF-DB / AffectNet / RAF-ML};
    \node[box, right=of data] (pre) {统一预处理\\resize / 标签映射 / 增强};
    \node[box, right=of pre] (model) {模型训练\\CNN / MiniResNet / RGBMiniResNet};
    \node[box, below=of model] (eval) {评估分析\\报告 / 混淆矩阵 / F1};
    \node[box, left=of eval] (ablation) {消融对比\\输入 / 策略 / 任务};
    \node[box, left=of ablation] (gui) {GUI 展示\\图片上传 / 摄像头识别};
    \draw[arrow] (data) -- (pre);
    \draw[arrow] (pre) -- (model);
    \draw[arrow] (model) -- (eval);
    \draw[arrow] (eval) -- (ablation);
    \draw[arrow] (ablation) -- (gui);
  \end{tikzpicture}
  \caption{项目整体实验流程}
  \label{fig:experiment_pipeline}
\end{figure}

从图~\ref{fig:experiment_pipeline} 可以看出，项目并不是单一模型训练流程，而是将多数据集输入、标签映射、模型比较、评估分析、消融实验和 GUI 展示串联成完整系统。主要训练实验的超参数设置见表~\ref{tab:experiment_hyperparams}，训练策略与权重文件见表~\ref{tab:experiment_weights}。

\begin{table}[H]
  \centering
  \caption{主要训练超参数设置}
  \label{tab:experiment_hyperparams}
  \scriptsize
  \setlength{\tabcolsep}{3pt}
  \renewcommand{\arraystretch}{1.12}
  \resizebox{\textwidth}{!}{%
  \begin{tabular}{llllllll}
    \toprule
    实验名称 & 数据集 & 输入形式 & Epochs & Batch Size & Learning Rate & Optimizer & Loss \\
    \midrule
    FER2013 MiniResNet & FER2013 & 灰度 48$\times$48 & 50 & 64 & 0.001 & Adam & CrossEntropyLoss \\
    RAF-DB RGBMiniResNet & RAF-DB Basic & RGB 96$\times$96 & 50 & 64 & 0.001 & Adam & CrossEntropyLoss \\
    AffectNet RGB Full & AffectNet YOLO Format & RGB 96$\times$96 & 40 & 128 & 0.001 & Adam & CrossEntropyLoss \\
    AffectNet 预训练 + RAF-DB 微调 & AffectNet $\rightarrow$ RAF-DB Basic & RGB 96$\times$96 & 50 & 64 & 0.0001 & Adam & CrossEntropyLoss \\
    RAF-DB Compound & RAF-DB Compound & RGB 96$\times$96 & 40 & 64 & 0.0005 & Adam & 加权 CrossEntropyLoss \\
    RAF-ML 多标签 & RAF-ML & RGB 96$\times$96 & 50 & 64 & 0.0005 & Adam & BCEWithLogitsLoss \\
    \bottomrule
  \end{tabular}
  }
\end{table}

\begin{table}[H]
  \centering
  \caption{训练策略与最优权重文件}
  \label{tab:experiment_weights}
  \small
  \setlength{\tabcolsep}{4pt}
  \renewcommand{\arraystretch}{1.15}
  \begin{tabularx}{\textwidth}{p{0.22\textwidth}p{0.10\textwidth}p{0.30\textwidth}X}
    \toprule
    实验名称 & 数据增强 & 正则化/备注 & 最优权重 \\
    \midrule
    FER2013 MiniResNet & 是 & 训练日志记录 50 个 epoch，使用脚本默认 Adam 设置。 & \path{checkpoints/best_mini_resnet.pth}\\
    RAF-DB RGBMiniResNet & 是 & 训练日志记录 50 个 epoch，使用 ReduceLROnPlateau 学习率调度。 & \path{checkpoints/best_rafdb_rgb_mini_resnet.pth}\\
    AffectNet RGB Full & 是 & 使用全量有效样本训练，日志记录 40 个 epoch。 & \path{checkpoints/best_affectnet_rgb_full_40epoch.pth}\\
    AffectNet 预训练 + RAF-DB 微调 & 是 & 加载 AffectNet RGB 权重初始化，在 RAF-DB Basic 上微调。 & \path{checkpoints/best_affectnet_pretrained_rafdb_rgb_mini_resnet.pth}\\
    RAF-DB Compound & 是 & 最佳模型为 AffectNet 预训练版本，日志记录 40 个 epoch，使用类别权重。 & \path{checkpoints/best_compound_rgb_affectnet_pretrained.pth}\\
    RAF-ML 多标签 & 是 & 最佳模型为 AffectNet 预训练版本，使用 sigmoid 阈值 0.5。 & \path{checkpoints/best_rafml_affectnet_pretrained.pth}\\
    \bottomrule
  \end{tabularx}
\end{table}

表~\ref{tab:experiment_hyperparams} 和表~\ref{tab:experiment_weights} 中的参数优先依据训练日志 CSV 的 epoch 记录、README 中的正式实验命令以及训练脚本默认值整理。与脚本默认轮数不同的实验以日志实际记录为准：RAF-DB RGB、AffectNet 预训练微调和 RAF-ML 均记录 50 个 epoch，Compound 最佳模型记录 40 个 epoch。各阶段实验安排如下：

\begin{enumerate}
  \item \textbf{FER2013 基础灰度模型训练：}训练 SimpleCNN、ImprovedCNN 和 MiniResNet，保存最优权重、训练日志和曲线。
  \item \textbf{CK+ 外部泛化测试：}使用 FER2013 训练的 MiniResNet 测试 CK+，删除 Contempt 后统一为七类。
  \item \textbf{RAF-DB 灰度外部测试：}将 RAF-DB RGB 图像转灰度后输入 FER2013 MiniResNet，分析真实场景跨数据集泛化能力。
  \item \textbf{RAF-DB RGB 从头训练：}使用 RAF-DB RGB 数据训练 RGBMiniResNet，输入为 96$\times$96 RGB 图像。
  \item \textbf{AffectNet RGB Full 训练：}使用 AffectNet YOLO Format 有效样本训练 7 类 RGB 模型。
  \item \textbf{AffectNet 预训练 + RAF-DB 微调：}先在 AffectNet 上预训练，再在 RAF-DB 上微调，分析迁移学习效果。
  \item \textbf{RAF-DB Compound 复合表情识别：}使用 11 类复合表情标签训练 RGBMiniResNet，并加入类别权重、label smoothing、dropout 和 early stopping 等正则化策略。
  \item \textbf{RAF-ML 多标签混合表情识别：}使用 6 维 sigmoid 多标签输出训练模型，评价指标使用 Micro F1 和 Macro F1。
  \item \textbf{消融实验整理：}通过 \code{analyze\_ablation\_results.py} 汇总已有结果，生成消融对比表和图表。
\end{enumerate}

在多标签混合表情任务中，训练流程与普通 softmax 分类不同。模型输出 6 个 logits，训练时直接使用 BCEWithLogitsLoss；评估时先经过 sigmoid 得到每个情绪成分概率，再根据阈值计算 Micro F1 和 Macro F1。核心训练逻辑如 Listing~\ref{lst:rafml_train} 所示。

\begin{lstlisting}[language=Python, caption={RAF-ML 多标签训练核心逻辑}, label={lst:rafml_train}]
criterion = nn.BCEWithLogitsLoss()

for images, labels, _ in loader:
    images = images.to(device)
    labels = labels.to(device)

    optimizer.zero_grad()
    logits = model(images)
    loss = criterion(logits, labels)
    loss.backward()
    optimizer.step()

    probabilities = torch.sigmoid(logits.detach())
    preds = (probabilities >= threshold).to(torch.int64)
\end{lstlisting}

Listing~\ref{lst:rafml_train} 展示了 RAF-ML 多标签任务中 logits、标签和 BCEWithLogitsLoss 的基本关系。与单标签任务不同，预测阶段需要对每个标签独立计算 sigmoid 概率并按阈值判断是否激活。

\section{GUI 系统设计}

\subsection{GUI 总体界面}

项目实现了基于 PyQt5 的桌面端 GUI，系统标题为“人脸表情识别系统”。界面左侧为图片与摄像头预览区域，右侧依次放置模型与任务设置、模型说明、预测结果和概率显示模块。系统支持基本表情识别、复合表情识别和混合表情识别三类任务，并通过两级下拉框完成任务类型与具体模型的选择。GUI 主界面如图~\ref{fig:gui_main} 所示。

\begin{figure}[H]
  \centering
  \figplaceholder{\rule{0pt}{1.2em}}
  \caption{GUI 主界面}
  \label{fig:gui_main}
\end{figure}

从图~\ref{fig:gui_main} 可以看出，GUI 将模型选择、权重路径、人脸定位方式、模型说明、预测结果和概率分布集中在同一界面中。切换任务或模型时，系统会自动更新权重路径、输出类别数、输出方式、标签列表、概率条和模型说明，从而降低不同实验模型之间的使用成本。

\subsection{图片上传识别功能}

图片上传识别模式用于对本地单张图像进行表情分类。用户选择图片后，系统使用原始图片文件或原始 PIL 图像进行预处理和模型预测，不受 GUI 显示框缩放尺寸影响。预测完成后，界面会显示预测类别、置信度和各类别概率；当切换到复合表情或混合表情任务时，概率条数量也会随任务自动变化。图片上传识别效果如图~\ref{fig:gui_image_prediction} 所示。

\begin{figure}[H]
  \centering
  \figplaceholder{\rule{0pt}{1.2em}}
  \caption{图片上传识别示例}
  \label{fig:gui_image_prediction}
\end{figure}

图~\ref{fig:gui_image_prediction} 展示了单张图片输入后的预测结果显示方式。对于 softmax 单标签任务，界面突出显示最高概率类别；对于 RAF-ML sigmoid 多标签任务，界面会显示超过阈值的多个情绪成分。

\subsection{摄像头实时识别功能}

摄像头实时识别模式通过 OpenCV 打开默认摄像头，并使用 QTimer 定时刷新画面，避免阻塞 GUI 主线程。系统支持 YuNet、MediaPipe 和中心区域识别三种人脸定位方式，其中人脸定位方式仅在摄像头模式下可选。每次预测时，系统基于摄像头原始帧进行定位与裁剪：若选择检测器并成功定位人脸，则使用人脸框区域进行表情识别；若检测器不可用或未检测到人脸，则回退到中心区域识别。停止摄像头后，界面会保留最后一帧，便于继续进行单帧预测。摄像头实时识别效果如图~\ref{fig:gui_camera_examples} 所示。为保护隐私，图中摄像头画面的人脸区域已做轻微模糊处理，但仍保留系统界面、预测结果和概率显示效果。

\begin{figure}[H]
  \centering
  \begin{subfigure}{0.48\textwidth}
    \centering
    \figplaceholder{\rule{0pt}{1.2em}}
    \caption{示例一}
  \end{subfigure}
  \hfill
  \begin{subfigure}{0.48\textwidth}
    \centering
    \figplaceholder{\rule{0pt}{1.2em}}
    \caption{示例二}
  \end{subfigure}
  \caption{摄像头实时识别示例}
  \label{fig:gui_camera_examples}
\end{figure}

图~\ref{fig:gui_camera_examples} 展示了摄像头模式下的实时画面、预测结果和概率分布。为保证交互流畅，模型不会在每一帧重复加载，而是在权重或模型配置变化后按需重新加载，并以固定间隔进行推理。

\subsection{多任务模型切换与动态输出}

GUI 支持三种任务输出形式：基本表情识别使用 7 类 softmax 单标签输出，复合表情识别使用 11 类 softmax 单标签输出，混合表情识别使用 6 维 sigmoid 多标签输出。三类任务在模型选择、输出方式和概率条数量上的差异见表~\ref{tab:gui_task_logic}。

\begin{table}[H]
  \centering
  \caption{GUI 支持的任务类型与输出逻辑}
  \label{tab:gui_task_logic}
  \scriptsize
  \setlength{\tabcolsep}{3pt}
  \renewcommand{\arraystretch}{1.18}
  \begin{tabularx}{\textwidth}{p{0.14\textwidth}p{0.32\textwidth}p{0.16\textwidth}p{0.10\textwidth}X}
    \toprule
    任务类型 & 可选模型 & 输出方式 & 概率条数量 & 适用场景 \\
    \midrule
    基本表情识别 & AffectNet RGB Full、RAF-DB RGBMiniResNet、AffectNet 预训练 + RAF-DB 微调、FER2013 灰度 MiniResNet & 7 类 softmax 单标签 & 7 & 图片上传和摄像头实时识别。\\
    复合表情识别 & RAF-DB Compound 复合表情 MiniResNet & 11 类 softmax 单标签 & 11 & 固定复合表情类别分析。\\
    混合表情识别 & RAF-ML 多标签混合表情 MiniResNet & 6 维 sigmoid 多标签 & 6 & 多个情绪成分同时存在的混合表情分析。\\
    \bottomrule
  \end{tabularx}
\end{table}

表~\ref{tab:gui_task_logic} 说明，GUI 将不同任务统一到同一界面中，但不会混淆不同任务的输出语义。人脸定位方式仅在摄像头实时识别模式下可选；图片上传识别模式下该选项显示“当前模式不可用”，图片预测继续使用原始图片预处理流程。由于 RAF-ML 允许多个情绪成分同时存在，因此不能被当作普通 softmax 单标签分类任务处理。GUI 的预测函数根据当前模型配置中的 \code{output\_type} 动态选择 softmax 或 sigmoid，核心逻辑如 Listing~\ref{lst:gui_output_switch} 所示。

\begin{lstlisting}[language=Python, caption={GUI 预测中 softmax / sigmoid 动态切换核心逻辑}, label={lst:gui_output_switch}]
logits = model(input_tensor.to(device))
if logits.shape[1] != len(labels_cn):
    raise ValueError("模型输出维度与当前任务标签数不一致。")

if output_name == "softmax":
    probs = torch.softmax(logits, dim=1)[0].detach().cpu().numpy()
elif output_name == "sigmoid":
    probs = torch.sigmoid(logits)[0].detach().cpu().numpy()
else:
    raise ValueError("不支持的输出类型，请选择 softmax 或 sigmoid。")

probability_list = [float(prob * 100.0) for prob in probs]
predicted_index = int(np.argmax(probs))

if output_name == "sigmoid":
    threshold_percent = float(threshold) * 100.0
    result["active_labels"] = [
        {"index": i, "label_cn": labels_cn[i],
         "label_en": labels_en[i], "probability": p}
        for i, p in enumerate(probability_list)
        if p >= threshold_percent
    ]
\end{lstlisting}

Listing~\ref{lst:gui_output_switch} 说明 GUI 根据模型配置动态选择 softmax 或 sigmoid，从而保证基本表情、复合表情和混合表情三类任务在同一界面中使用正确的输出解释方式。这一设计避免了将 RAF-ML 多标签模型误当作普通七类分类模型，也使 Compound 11 类模型能够复用同一套界面输出逻辑。

\section{实验结果与分析}

\subsection{总体结果}

各阶段主要实验结果如表~\ref{tab:main_results} 所示。该表将外部泛化测试、RGB 同源训练、AffectNet 迁移、Compound 复合表情和 RAF-ML 多标签任务放在同一视角下比较，便于观察不同任务定义和数据来源对模型性能的影响。

{\scriptsize
\setlength{\tabcolsep}{2pt}
\renewcommand{\arraystretch}{1.15}
\begin{longtable}{p{0.18\textwidth}p{0.13\textwidth}p{0.15\textwidth}p{0.08\textwidth}p{0.08\textwidth}p{0.08\textwidth}p{0.18\textwidth}}
  \caption{主要实验结果汇总}
  \label{tab:main_results}\\
  \toprule
  实验名称 & 测试数据集 & 任务类型 & Accuracy & Micro F1 & Macro F1 & 说明 \\
  \midrule
  \endfirsthead
  \toprule
  实验名称 & 测试数据集 & 任务类型 & Accuracy & Micro F1 & Macro F1 & 说明 \\
  \midrule
  \endhead
  CK+ 外部测试 & CK+ & 7 类 softmax 单标签 & 0.7140 & -- & 0.5275 & FER2013 灰度 MiniResNet 跨数据集测试。\\
  RAF-DB 灰度外部测试 & RAF-DB test & 7 类 softmax 单标签 & 0.5684 & -- & 0.3921 & FER2013 灰度模型迁移到真实场景数据。\\
  RAF-DB RGB 从头训练 & RAF-DB test & 7 类 softmax 单标签 & 0.8380 & -- & 0.7591 & RGBMiniResNet 在同源真实场景数据上表现最好。\\
  AffectNet RGB Full & AffectNet test & 7 类 softmax 单标签 & 0.7594 & -- & 0.7491 & 真实网络图像训练，适合摄像头演示。\\
  AffectNet 预训练 + RAF-DB 微调 & RAF-DB test & 7 类 softmax 单标签 & 0.8181 & -- & 0.7302 & 迁移学习模型，接近 RAF-DB 从头训练。\\
  RAF-DB Compound & Compound test & 11 类 softmax 单标签 & 0.5846 & -- & 0.4492 & 复合表情固定类别识别，难度较高。\\
  RAF-ML 多标签混合表情 & RAF-ML test & 6 维 sigmoid 多标签 & -- & 0.7588 & 0.7352 & 多标签任务不使用普通单标签 Accuracy。\\
  \bottomrule
\end{longtable}
}

从表~\ref{tab:main_results} 可以看出，RAF-DB RGB 从头训练在基本七类表情识别中表现最好；AffectNet 预训练 + RAF-DB 微调接近但略低于同源训练；Compound 复合表情任务明显更难；RAF-ML 多标签任务使用 Micro F1 和 Macro F1 更适合评价。

\subsection{FER2013 灰度模型对比分析}

FER2013 阶段训练了 SimpleCNN、ImprovedCNN 和 MiniResNet 三个灰度模型，用于逐步比较基础 CNN、常规改进 CNN 和轻量残差网络的差异。三者输入均为 48$\times$48 灰度图，输出均为 7 类 softmax 单标签，但网络结构和正则化程度不同。SimpleCNN、ImprovedCNN 和 MiniResNet 的训练曲线分别见图~\ref{fig:simple_cnn_curves}、图~\ref{fig:improved_cnn_curves} 和图~\ref{fig:mini_resnet_curves}。

三种 FER2013 灰度模型在测试集上的主要结果见表~\ref{tab:fer2013_gray_compare}。该表汇总了分类报告中的总体准确率和宏平均指标，并补充了主要模型的参数量记录。

\begin{table}[H]
  \centering
  \caption{FER2013 灰度模型对比结果}
  \label{tab:fer2013_gray_compare}
  \small
  \setlength{\tabcolsep}{4pt}
  \renewcommand{\arraystretch}{1.15}
  \begin{tabularx}{\textwidth}{lccccp{0.13\textwidth}X}
    \toprule
    模型 & Accuracy & Macro Precision & Macro Recall & Macro F1 & 参数量 & 说明 \\
    \midrule
    SimpleCNN & 0.5074 & 0.4918 & 0.4596 & 0.4621 & -- & 基础 CNN，对训练评估流程进行验证。\\
    ImprovedCNN & 0.6431 & 0.6389 & 0.5796 & 0.5909 & 1,320,871 & 加入 BatchNorm、Dropout 和更深卷积结构。\\
    MiniResNet & 0.6520 & 0.6419 & 0.6304 & 0.6325 & 695,655 & 轻量残差网络，作为灰度主实验模型。\\
    \bottomrule
  \end{tabularx}
\end{table}

表~\ref{tab:fer2013_gray_compare} 表明，SimpleCNN 作为基础模型效果最低；ImprovedCNN 通过更深网络、BatchNorm 和 Dropout 明显提升了测试集表现；MiniResNet 的参数量小于 ImprovedCNN，但 Macro F1 更高，说明残差连接在 FER2013 灰度表情识别中更有效。因此后续 CK+ 外部测试和 RAF-DB 灰度外部测试均选用 MiniResNet 作为灰度代表模型。

\begin{figure}[H]
  \centering
  \begin{subfigure}{0.48\textwidth}
    \centering
    \figplaceholder{\rule{0pt}{1.2em}}
    \caption{SimpleCNN Accuracy 曲线}
  \end{subfigure}
  \hfill
  \begin{subfigure}{0.48\textwidth}
    \centering
    \figplaceholder{\rule{0pt}{1.2em}}
    \caption{SimpleCNN Loss 曲线}
  \end{subfigure}
  \caption{SimpleCNN 在 FER2013 上的训练曲线}
  \label{fig:simple_cnn_curves}
\end{figure}

\begin{figure}[H]
  \centering
  \begin{subfigure}{0.48\textwidth}
    \centering
    \figplaceholder{\rule{0pt}{1.2em}}
    \caption{ImprovedCNN Accuracy 曲线}
  \end{subfigure}
  \hfill
  \begin{subfigure}{0.48\textwidth}
    \centering
    \figplaceholder{\rule{0pt}{1.2em}}
    \caption{ImprovedCNN Loss 曲线}
  \end{subfigure}
  \caption{ImprovedCNN 在 FER2013 上的训练曲线}
  \label{fig:improved_cnn_curves}
\end{figure}

\begin{figure}[H]
  \centering
  \begin{subfigure}{0.48\textwidth}
    \centering
    \figplaceholder{\rule{0pt}{1.2em}}
    \caption{MiniResNet Accuracy 曲线}
  \end{subfigure}
  \hfill
  \begin{subfigure}{0.48\textwidth}
    \centering
    \figplaceholder{\rule{0pt}{1.2em}}
    \caption{MiniResNet Loss 曲线}
  \end{subfigure}
  \caption{MiniResNet 在 FER2013 上的训练曲线}
  \label{fig:mini_resnet_curves}
\end{figure}

从三组 FER2013 曲线可以看出，模型结构从 SimpleCNN 到 ImprovedCNN 再到 MiniResNet 后，训练过程逐渐具备更强的特征表达能力。SimpleCNN 更适合作为流程验证基线；ImprovedCNN 通过 BatchNorm 和 Dropout 改善训练稳定性；MiniResNet 通过残差连接进一步增强灰度表情特征提取能力，因此被选为 FER2013 主实验和外部泛化测试的代表模型。

\subsection{基本表情识别结果分析}

基本七类表情识别中，RAF-DB RGBMiniResNet 在 RAF-DB test 上取得 Accuracy 0.8380 和 Macro F1 0.7591，效果明显优于 RAF-DB 灰度外部测试结果。该现象说明，真实场景彩色图像中的颜色、光照和纹理信息对表情识别有帮助，同时同源训练和测试能够减小数据分布差异。RAF-DB RGBMiniResNet 的训练曲线如图~\ref{fig:rafdb_curves_in_results} 所示。

\begin{figure}[H]
  \centering
  \begin{subfigure}{0.48\textwidth}
    \centering
    \figplaceholder{\rule{0pt}{1.2em}}
    \caption{Accuracy 曲线}
  \end{subfigure}
  \hfill
  \begin{subfigure}{0.48\textwidth}
    \centering
    \figplaceholder{\rule{0pt}{1.2em}}
    \caption{Loss 曲线}
  \end{subfigure}
  \caption{RAF-DB RGBMiniResNet 训练曲线}
  \label{fig:rafdb_curves_in_results}
\end{figure}

从 RAF-DB RGB 训练曲线可以看到，训练集 accuracy 持续上升，验证集 accuracy 也逐步提高但存在一定波动，说明模型整体已经收敛。训练 loss 持续下降，而验证 loss 在后期趋于平稳并伴随波动，说明模型对 RAF-DB Basic 具有较好的拟合能力，但也存在一定过拟合倾向。

AffectNet RGB Full 模型在 AffectNet test 上取得 Accuracy 0.7594 和 Macro F1 0.7491，说明在复杂网络图像场景中也具有较稳定的识别能力，其训练曲线如图~\ref{fig:affectnet_full_curves_in_results} 所示。AffectNet 预训练 + RAF-DB 微调模型在 RAF-DB test 上取得 Accuracy 0.8181 和 Macro F1 0.7302，接近 RAF-DB 从头训练模型，但仍略低；对应微调曲线见图~\ref{fig:affectnet_pretrained_rafdb_curves}。

\begin{figure}[H]
  \centering
  \begin{subfigure}{0.48\textwidth}
    \centering
    \figplaceholder{\rule{0pt}{1.2em}}
    \caption{Accuracy 曲线}
  \end{subfigure}
  \hfill
  \begin{subfigure}{0.48\textwidth}
    \centering
    \figplaceholder{\rule{0pt}{1.2em}}
    \caption{Loss 曲线}
  \end{subfigure}
  \caption{AffectNet RGB Full MiniResNet 训练曲线}
  \label{fig:affectnet_full_curves_in_results}
\end{figure}

AffectNet RGB Full 的训练曲线整体更平滑，验证性能随训练逐步提升，说明模型能够从更大规模、更复杂的真实场景图像中学习到较稳定的 RGB 表情特征。由于 AffectNet 图像来源复杂且标注噪声相对更高，最终性能低于 RAF-DB 同源测试并不意外，但该模型在摄像头演示中表现更稳定，因此被 GUI 作为默认推荐模型。

\begin{figure}[H]
  \centering
  \begin{subfigure}{0.48\textwidth}
    \centering
    \figplaceholder{\rule{0pt}{1.2em}}
    \caption{Accuracy 曲线}
  \end{subfigure}
  \hfill
  \begin{subfigure}{0.48\textwidth}
    \centering
    \figplaceholder{\rule{0pt}{1.2em}}
    \caption{Loss 曲线}
  \end{subfigure}
  \caption{AffectNet 预训练 + RAF-DB 微调训练曲线}
  \label{fig:affectnet_pretrained_rafdb_curves}
\end{figure}

AffectNet 预训练 + RAF-DB 微调曲线显示，预训练初始化使模型在 RAF-DB 上能够较快达到较高性能，说明 AffectNet 中学习到的真实场景表情特征具有一定迁移价值。但该模型最终仍略低于 RAF-DB 从头训练模型，可能与 AffectNet 和 RAF-DB 在图像来源、标签映射、类别比例和标注风格上的差异有关。AffectNet RGB Full 与 RAF-DB RGB 的混淆矩阵分别见图~\ref{fig:affectnet_confusion_in_results} 和图~\ref{fig:rafdb_confusion_in_results}。

\begin{figure}[H]
  \centering
  \figplaceholder{\rule{0pt}{1.2em}}
  \caption{AffectNet RGB Full MiniResNet 混淆矩阵}
  \label{fig:affectnet_confusion_in_results}
\end{figure}

\begin{figure}[H]
  \centering
  \figplaceholder{\rule{0pt}{1.2em}}
  \caption{RAF-DB RGBMiniResNet 混淆矩阵}
  \label{fig:rafdb_confusion_in_results}
\end{figure}

从 RAF-DB RGB 混淆矩阵可以观察基本七类表情之间的错误分布。通常 Happy、Surprise 等表情具有更明显的视觉特征，模型更容易识别；Fear、Sad、Neutral 等类别在真实场景中表情强度较弱，且局部面部变化更接近，因此更容易互相混淆。该现象也解释了为什么只看 Accuracy 不够，还需要结合 Macro F1 和每类 Precision/Recall 分析小类表现。

RAF-DB RGBMiniResNet 的每类指标见表~\ref{tab:rafdb_per_class}。该表进一步说明，类别样本数和视觉区分度都会影响模型表现。

\begin{table}[H]
  \centering
  \caption{RAF-DB RGBMiniResNet 每类识别指标}
  \label{tab:rafdb_per_class}
  \small
  \setlength{\tabcolsep}{6pt}
  \begin{tabular}{lrrrr}
    \toprule
    类别 & Precision & Recall & F1 & Support \\
    \midrule
    Angry & 0.7278 & 0.8086 & 0.7661 & 162 \\
    Disgust & 0.6210 & 0.4813 & 0.5423 & 160 \\
    Fear & 0.5974 & 0.6216 & 0.6093 & 74 \\
    Happy & 0.9403 & 0.9038 & 0.9217 & 1185 \\
    Sad & 0.7971 & 0.7971 & 0.7971 & 478 \\
    Surprise & 0.8416 & 0.8723 & 0.8567 & 329 \\
    Neutral & 0.7929 & 0.8500 & 0.8204 & 680 \\
    \bottomrule
  \end{tabular}
\end{table}

从表~\ref{tab:rafdb_per_class} 可以看出，Happy 和 Surprise 的 F1 分别达到 0.9217 和 0.8567，表现相对较好；Disgust 和 Fear 的 F1 分别为 0.5423 和 0.6093，相对较低。Happy、Surprise 通常具有较明显的嘴部或眼部变化，而 Fear、Disgust 等类别样本更少，且与 Sad、Neutral 等类别在真实图像中可能存在局部表情强度弱、视觉差异细的问题，因此更容易混淆。

\subsection{外部泛化测试分析}

外部泛化测试用于观察模型在训练数据之外的数据集上的表现。FER2013 灰度 MiniResNet 分别在 CK+ 和 RAF-DB 灰度外部测试上进行评估，结果见表~\ref{tab:external_generalization}。

\begin{table}[H]
  \centering
  \caption{外部泛化测试结果对比}
  \label{tab:external_generalization}
  \small
  \setlength{\tabcolsep}{4pt}
  \renewcommand{\arraystretch}{1.15}
  \begin{tabularx}{\textwidth}{p{0.14\textwidth}p{0.13\textwidth}ccp{0.28\textwidth}X}
    \toprule
    测试数据集 & 输入形式 & Accuracy & Macro F1 & 数据特点 & 结论 \\
    \midrule
    CK+ & 灰度 48$\times$48 & 0.7140 & 0.5275 & 实验室环境采集，表情更清晰、更夸张，背景和姿态变化相对受控。 & 灰度 MiniResNet 能取得较高 Accuracy，但 Macro F1 仍受类别不平衡影响。\\
    RAF-DB test & 灰度 48$\times$48 & 0.5684 & 0.3921 & 真实网络场景图像，存在姿态、光照、遮挡、图像质量、年龄和场景差异。 & 跨到真实场景后性能明显下降，说明灰度 FER2013 模型泛化能力有限。\\
    \bottomrule
  \end{tabularx}
\end{table}

从表~\ref{tab:external_generalization} 可以看出，同一个 FER2013 灰度 MiniResNet 在 CK+ 和 RAF-DB 上的表现差异明显。CK+ 的采集环境更受控，表情通常更清晰、更夸张，因此模型即使没有在 CK+ 上训练，也能取得 0.7140 的 Accuracy。相比之下，RAF-DB 来源于真实场景，图像中存在更复杂的姿态、光照、遮挡、图像清晰度和背景变化，将彩色图像转为灰度后还会丢失颜色与纹理信息，因此 Accuracy 和 Macro F1 均明显下降。

两组外部测试还说明，只在单一数据集内部报告测试结果并不足以全面反映模型真实泛化能力。CK+ 与 RAF-DB 的差异表明，数据采集环境和图像风格会显著影响表情识别模型性能。此外，两组实验中的 Macro F1 均低于 Accuracy，尤其 RAF-DB 灰度外部测试 Macro F1 只有 0.3921，说明跨数据集时小类表现和类别不平衡问题更加突出。因此，本项目后续引入 RAF-DB RGB 同源训练、AffectNet 预训练以及消融实验，用于更全面地分析真实场景下的表情识别能力。

\subsection{复合表情识别分析}

RAF-DB Compound 复合表情识别取得 Accuracy 0.5846 和 Macro F1 0.4492，明显低于基本七类表情识别。该任务包含 11 类更细粒度的复合表情，类别之间语义接近，例如“悲伤-恐惧”和“悲伤-惊讶”之间边界较模糊。同时 Compound 数据量少于 RAF-DB Basic，类别不平衡问题更明显，模型容易出现过拟合。Compound 最佳模型训练曲线如图~\ref{fig:compound_curves_in_results} 所示。

\begin{figure}[H]
  \centering
  \begin{subfigure}{0.48\textwidth}
    \centering
    \figplaceholder{\rule{0pt}{1.2em}}
    \caption{Accuracy 曲线}
  \end{subfigure}
  \hfill
  \begin{subfigure}{0.48\textwidth}
    \centering
    \figplaceholder{\rule{0pt}{1.2em}}
    \caption{Loss 曲线}
  \end{subfigure}
  \caption{RAF-DB Compound 训练曲线}
  \label{fig:compound_curves_in_results}
\end{figure}

图~\ref{fig:compound_curves_in_results} 使用的是 \code{compound\_rgb\_affectnet\_pretrained} 最佳模型对应曲线。曲线中训练 accuracy 明显高于验证或测试 accuracy，loss 差距也较明显，说明该任务过拟合更明显。主要原因是 Compound 数据量较少，11 个固定复合类别之间语义细粒度且类别边界接近，模型更容易记住训练集特征，而难以在测试集上稳定区分相近复合情绪。Compound 混淆矩阵如图~\ref{fig:compound_confusion_in_results} 所示。

\begin{figure}[H]
  \centering
  \figplaceholder{\rule{0pt}{1.2em}}
  \caption{RAF-DB Compound 混淆矩阵}
  \label{fig:compound_confusion_in_results}
\end{figure}

Compound 混淆矩阵进一步说明，固定复合表情类别之间的混淆明显高于基本七类任务。复合标签常常共享一个基础情绪成分，例如多个类别都包含 Sadly、Fearfully 或 Angrily，因此模型需要区分更细微的第二情绪成分。由于样本数量较少且类别语义接近，模型即使在训练集上达到较高准确率，也容易在测试集上出现过拟合和跨类别混淆。

RAF-DB Compound 每类 F1 指标见表~\ref{tab:compound_per_class_f1}。与基本七类任务相比，复合类别之间的性能差异更加明显。

\begin{table}[H]
  \centering
  \caption{RAF-DB Compound 各复合表情类别 F1 指标}
  \label{tab:compound_per_class_f1}
  \small
  \setlength{\tabcolsep}{5pt}
  \renewcommand{\arraystretch}{1.12}
  \begin{tabularx}{\textwidth}{p{0.19\textwidth}p{0.28\textwidth}ccc}
    \toprule
    中文类别 & 英文类别 & F1 & Support & 说明 \\
    \midrule
    开心-惊讶 & Happily Surprised & 0.7114 & 135 & 表现较好 \\
    开心-厌恶 & Happily Disgusted & 0.4719 & 47 & 中等 \\
    悲伤-恐惧 & Sadly Fearful & 0.1860 & 22 & 较难 \\
    悲伤-愤怒 & Sadly Angry & 0.2769 & 33 & 较难 \\
    悲伤-惊讶 & Sadly Surprised & 0.2667 & 18 & 较难 \\
    悲伤-厌恶 & Sadly Disgusted & 0.6301 & 141 & 表现较好 \\
    恐惧-愤怒 & Fearfully Angry & 0.5455 & 33 & 中等 \\
    恐惧-惊讶 & Fearfully Surprised & 0.6667 & 116 & 表现较好 \\
    愤怒-惊讶 & Angrily Surprised & 0.3288 & 38 & 较难 \\
    愤怒-厌恶 & Angrily Disgusted & 0.6685 & 174 & 表现较好 \\
    厌恶-惊讶 & Disgustedly Surprised & 0.1887 & 35 & 较难 \\
    \bottomrule
  \end{tabularx}
\end{table}

表~\ref{tab:compound_per_class_f1} 显示，Happily Surprised、Angrily Disgusted、Fearfully Surprised 等类别较容易识别，而 Sadly Fearful、Sadly Surprised、Disgustedly Surprised 等类别较难。原因在于复合表情中的第二情绪成分往往不够明显，并且部分类别样本数较少，模型需要在非常接近的面部细节之间做出区分，因此测试集性能容易下降。

\subsection{RAF-ML 多标签结果分析}

RAF-ML 多标签混合表情识别取得 Micro F1 0.7588 和 Macro F1 0.7352，明显高于 RAF-DB Compound 11 类固定复合表情结果。原因在于 RAF-ML 将混合表情拆解为多个基础情绪成分，每个维度独立建模，更符合真实场景中多种情绪同时存在的特点。RAF-ML 多标签模型的 F1 和 loss 曲线如图~\ref{fig:rafml_f1_curve_in_results} 所示。

\begin{figure}[H]
  \centering
  \begin{subfigure}{0.48\textwidth}
    \centering
    \figplaceholder{\rule{0pt}{1.2em}}
    \caption{训练集与测试集 Micro/Macro F1 曲线}
  \end{subfigure}
  \hfill
  \begin{subfigure}{0.48\textwidth}
    \centering
    \figplaceholder{\rule{0pt}{1.2em}}
    \caption{Loss 曲线}
  \end{subfigure}
  \caption{RAF-ML 多标签模型训练曲线}
  \label{fig:rafml_f1_curve_in_results}
\end{figure}

RAF-ML 的训练集与测试集 Micro/Macro F1 曲线整体逐渐上升并趋于稳定，说明多标签训练过程有效。测试 Macro F1 通常低于 Micro F1，说明不同标签之间存在难度差异和类别分布不均衡；loss 曲线则用于观察多标签 BCE 训练是否稳定收敛。因此 RAF-ML 更适合用 Micro F1、Macro F1 和每个标签的 Precision/Recall 共同分析，而不是使用单标签 Accuracy。

RAF-ML 每个情绪标签的 Precision、Recall 和 F1 见表~\ref{tab:rafml_per_label}。该表能够更细致地观察不同情绪成分的识别难度。

\begin{table}[H]
  \centering
  \caption{RAF-ML 多标签模型每标签识别指标}
  \label{tab:rafml_per_label}
  \small
  \setlength{\tabcolsep}{6pt}
  \begin{tabular}{lrrrr}
    \toprule
    标签 & Precision & Recall & F1 & Support \\
    \midrule
    Surprise & 0.8469 & 0.8137 & 0.8300 & 408 \\
    Fear & 0.7514 & 0.5603 & 0.6420 & 232 \\
    Disgust & 0.8459 & 0.8553 & 0.8506 & 539 \\
    Happiness & 0.6462 & 0.6512 & 0.6486 & 258 \\
    Sadness & 0.7081 & 0.6647 & 0.6857 & 343 \\
    Anger & 0.7565 & 0.7526 & 0.7546 & 384 \\
    \bottomrule
  \end{tabular}
\end{table}

从表~\ref{tab:rafml_per_label} 可以看出，Disgust、Surprise 和 Anger 表现较好，其中 Disgust 的 F1 为 0.8506，Surprise 的 F1 为 0.8300。Fear 和 Happiness 相对较低，Fear 的 Recall 只有 0.5603，可能与正样本数量较少、表情强度弱以及与 Surprise、Sadness 存在视觉相似有关；Happiness 在多标签任务中可能与 Surprise 等标签共现，其边界不如单标签 Happy 明确，因此 F1 低于 RAF-DB Basic 中的 Happy 类。

与 Compound 把“开心-惊讶”等组合视为一个互斥类别不同，RAF-ML 允许 Happiness 和 Surprise 同时为正标签。这样的任务定义降低了固定组合类别之间的硬边界问题，也更符合真实表情的连续性和混合性。因此 RAF-ML 在 Macro F1 上更稳定，但它没有 Neutral 类，也不能用普通单标签 Accuracy 评价。

\section{消融实验与对比分析}

为进一步分析不同实验因素对表情识别性能的影响，项目设计了输入形式、训练策略和任务形式三类消融对比。消融分析脚本中的柱状图生成核心逻辑见 Listing~\ref{lst:ablation_plot}。

\begin{lstlisting}[language=Python, caption={消融分析脚本生成柱状图的核心逻辑}, label={lst:ablation_plot}]
def plot_grouped_bar(labels, metric_names, values_by_metric, title, output_path):
    x_positions = list(range(len(labels)))
    width = 0.32
    fig, ax = plt.subplots(figsize=(8.2, 5.0), dpi=150)

    for metric_index, (metric_name, values) in enumerate(zip(metric_names, values_by_metric)):
        offset = (metric_index - (len(metric_names) - 1) / 2) * width
        bars = ax.bar(
            [x + offset for x in x_positions],
            values,
            width,
            label=metric_name,
        )
        add_bar_labels(ax, bars)

    ax.set_title(title, fontsize=14, fontweight="bold")
    ax.set_ylim(0, 1.0)
    ax.set_xticks(x_positions)
    ax.set_xticklabels(labels, rotation=12, ha="right")
    ax.legend()
    fig.tight_layout()
    fig.savefig(output_path)
\end{lstlisting}

Listing~\ref{lst:ablation_plot} 使用分组柱状图比较不同实验设置下的指标差异，并在柱状图上标注数值，便于直观说明性能变化。

\subsection{输入形式消融：灰度图 vs RGB 彩色图}

输入形式消融比较 FER2013 灰度 MiniResNet 在 RAF-DB 外部测试上的结果与 RAF-DB RGBMiniResNet 在同源 test set 上的结果，具体数值见表~\ref{tab:input_ablation}，对应柱状图见图~\ref{fig:input_ablation}。

\begin{table}[H]
  \centering
  \caption{输入形式消融结果}
  \label{tab:input_ablation}
  \begin{tabular}{lccc}
    \toprule
    模型 & 输入形式 & Accuracy & Macro F1 \\
    \midrule
    FER2013 灰度 MiniResNet 外部测试 & 灰度 48$\times$48 & 0.5684 & 0.3921 \\
    RAF-DB RGBMiniResNet & RGB 96$\times$96 & 0.8380 & 0.7591 \\
    \bottomrule
  \end{tabular}
\end{table}

\begin{figure}[H]
  \centering
  \figplaceholder{\rule{0pt}{1.2em}}
  \caption{输入形式消融：灰度图 vs RGB 彩色图}
  \label{fig:input_ablation}
\end{figure}

输入形式消融结果表明，FER2013 灰度 MiniResNet 在 RAF-DB 外部测试上的表现明显低于 RAF-DB RGBMiniResNet。前者同时受到灰度输入、低分辨率和跨数据集域差异影响；后者使用 RAF-DB 同源彩色真实场景图像训练，保留了更多颜色、纹理和光照信息，因此 Accuracy 和 Macro F1 均显著提升。

\subsection{训练策略消融：从头训练 vs AffectNet 预训练微调}

训练策略消融比较 RAF-DB RGB 从头训练模型与 AffectNet 预训练 + RAF-DB 微调模型在 RAF-DB test 上的表现，结果见表~\ref{tab:training_strategy_ablation}，可视化柱状图见图~\ref{fig:training_strategy_ablation}。

\begin{table}[H]
  \centering
  \caption{训练策略消融结果}
  \label{tab:training_strategy_ablation}
  \begin{tabular}{lccc}
    \toprule
    模型 & 训练策略 & Accuracy & Macro F1 \\
    \midrule
    RAF-DB RGBMiniResNet & 从头训练 & 0.8380 & 0.7591 \\
    AffectNet 预训练 + RAF-DB 微调 & 预训练微调 & 0.8181 & 0.7302 \\
    \bottomrule
  \end{tabular}
\end{table}

\begin{figure}[H]
  \centering
  \figplaceholder{\rule{0pt}{1.2em}}
  \caption{训练策略消融：从头训练 vs AffectNet 预训练微调}
  \label{fig:training_strategy_ablation}
\end{figure}

训练策略消融结果显示，AffectNet 预训练 + RAF-DB 微调取得了接近 RAF-DB 从头训练的性能，但在 RAF-DB test 上仍略低于同源数据从头训练模型。这说明 AffectNet 预训练可以提供有用的表情特征初始化，但 AffectNet 与 RAF-DB 在图像来源、标签映射、类别比例和采集风格上存在差异，可能限制了迁移收益。

\subsection{任务形式对比：基本表情 vs 复合表情 vs 混合表情}

任务形式对比用于说明基本七类 softmax 单标签、Compound 11 类 softmax 单标签和 RAF-ML 6 维 sigmoid 多标签任务的难度差异，结果见表~\ref{tab:task_comparison}，柱状图见图~\ref{fig:task_comparison}。

\begin{table}[H]
  \centering
  \caption{任务形式对比结果}
  \label{tab:task_comparison}
  \begin{tabular}{llcc}
    \toprule
    任务 & 输出方式 & 主指标 & Macro F1 \\
    \midrule
    基本七类表情识别 & 7 类 softmax 单标签 & Accuracy / Macro F1 & 0.7591 \\
    Compound 11 类复合表情识别 & 11 类 softmax 单标签 & Accuracy / Macro F1 & 0.4492 \\
    RAF-ML 6 维混合表情识别 & 6 维 sigmoid 多标签 & Micro F1 / Macro F1 & 0.7352 \\
    \bottomrule
  \end{tabular}
\end{table}

\begin{figure}[H]
  \centering
  \figplaceholder{\rule{0pt}{1.2em}}
  \caption{任务形式对比：基本表情 vs 复合表情 vs 混合表情}
  \label{fig:task_comparison}
\end{figure}

任务形式对比说明，基本七类表情识别相对更稳定。Compound 11 类固定复合表情识别难度更高，主要原因是样本数量较少、类别更细、语义更接近。RAF-ML 多标签混合表情识别 Macro F1 较高，说明将混合情绪拆解为多个基础情绪成分，更适合表达真实场景中多个情绪同时存在的情况。

\section{可视化分析}

项目生成了训练 Loss 曲线、训练 Accuracy 或 F1 曲线、混淆矩阵、分类报告、类别统计、正确样例、错误样例和消融实验柱状图。主要可视化结果按照实验对象分布在对应分析位置：RAF-DB RGB、AffectNet RGB Full、RAF-DB Compound 和 RAF-ML 的训练曲线分别放在各自结果分析小节；RAF-DB RGB、AffectNet RGB Full 和 Compound 的混淆矩阵也放在对应模型分析下；三张消融柱状图保留在“消融实验与对比分析”章节。

\subsection{训练曲线分析}

训练曲线用于观察模型是否收敛、是否出现过拟合，以及训练集和验证/测试集之间是否存在明显差距。对于 RAF-DB RGBMiniResNet，训练 accuracy 持续上升而验证 accuracy 有一定波动，说明模型能够学习真实场景表情特征，但后期仍存在过拟合趋势。AffectNet RGB Full 曲线整体较平滑，说明更大规模真实场景数据有助于学习稳定特征。Compound 曲线中训练 accuracy 明显高于测试 accuracy，体现出复合表情小样本、细粒度和类别不平衡带来的过拟合。RAF-ML 曲线使用 Micro/Macro F1 而不是单标签 accuracy，更适合观察多标签任务中各情绪维度的整体学习趋势。

\subsection{混淆矩阵分析}

混淆矩阵用于观察类别之间的误判关系。基本七类表情中，Happy 和 Surprise 通常具有较明显的面部变化，因此更容易形成清晰的对角线；Fear、Sad、Neutral 等类别表情强度较弱，在真实场景中更容易互相混淆。Compound 任务的混淆矩阵则更复杂，因为多个复合类别共享同一个基础情绪成分，例如 Sadly、Fearfully 或 Angrily，模型需要进一步区分第二情绪成分。相比之下，RAF-ML 是多标签 sigmoid 任务，不适合直接使用普通单标签混淆矩阵，更适合结合每个标签的 Precision、Recall、F1 和阈值分析进行解释。

\subsection{正确样例与错误样例分析}

项目在评估阶段保存了 correct\_samples 和 wrong\_samples 相关结果，用于观察模型在典型表情和困难样本上的表现。正确样例通常具有较明显的表情特征、人脸区域清晰、光照稳定、遮挡较少、类别边界清楚等特点。例如 Happy 和 Surprise 往往伴随明显的嘴角、眼部或眉部变化，因此更容易被模型稳定识别。

错误样例的来源更加复杂。第一，弱表情样本中面部动作幅度较小，Fear、Sad、Neutral 等类别在视觉上容易接近，模型可能难以捕捉细微差异。第二，侧脸、姿态偏转、光照过暗或过曝、人脸遮挡、图像模糊等因素会降低有效面部特征质量。第三，真实场景图像中标签本身具有一定主观性，同一张人脸可能同时带有多种情绪成分，使单标签模型被迫选择一个主类别。第四，在 Compound 任务中，多个类别共享相同基础情绪成分，例如 Sadly、Fearfully、Angrily 等，模型不仅要识别基础情绪，还要判断第二情绪成分，因而更容易在语义接近的复合类别之间混淆。

对于 RAF-ML 多标签任务，错误样例不一定意味着整张图像完全识别失败。有时模型已经给出某个情绪成分较高的概率，但该概率未超过固定阈值 0.5，因此在二值化后被计为漏检；也有可能多个相关情绪同时被激活，导致 Precision 和 Recall 随阈值变化而明显波动。因此，RAF-ML 后续可以尝试动态阈值、类别自适应阈值或标签相关性建模，进一步分析 Happiness--Surprise、Fear--Sadness 等情绪成分之间的共现关系。

总体而言，正确和错误样例分析说明，表情识别不仅是分类器结构问题，还与数据质量、人脸定位、表情强度、标签主观性和任务定义密切相关。将分类报告、混淆矩阵、类别统计与样例可视化结合起来，能够比单独观察 Accuracy 更准确地解释模型成功和失败的原因。

\section{创新点与项目特色}

本项目相比普通图像分类任务，不仅完成单一模型训练，还从实验设计、任务扩展和系统实现三个层面进行了综合探索。

\textbf{实验设计层面：}项目使用 FER2013、CK+、RAF-DB、AffectNet、RAF-DB Compound 和 RAF-ML 等多个数据集，覆盖基础训练、外部泛化测试、真实场景 RGB 训练和扩展任务实验。实验中对比了灰度输入与 RGB 输入、FER2013 外部测试与 RAF-DB 同源训练、RAF-DB 从头训练与 AffectNet 预训练微调，并通过消融分析脚本整理输入形式、训练策略和任务形式对模型性能的影响。

\textbf{任务扩展层面：}项目从基本七类表情识别进一步扩展到 RAF-DB Compound 11 类复合表情识别和 RAF-ML 6 维多标签混合表情识别。基本表情和复合表情采用 softmax 单标签分类，RAF-ML 采用 sigmoid 多标签输出，使系统能够同时处理互斥类别分类和多情绪成分识别两类任务。这样的设计体现了真实表情从单一基本情绪到复杂混合情绪的任务扩展过程。

\textbf{系统实现层面：}项目搭建了 PyQt5 交互式 GUI，支持图片上传识别、摄像头实时识别、两级模型选择、动态概率条和模型说明。摄像头端支持 YuNet、MediaPipe 和中心区域识别三种人脸定位方式；GUI 根据不同任务自动切换标签列表、输出方式和概率条数量，从而避免将 softmax 单标签任务和 sigmoid 多标签任务混用。整体系统将离线训练模型转化为可交互演示应用，增强了实验结果的可展示性和可验证性。

\section{不足与后续工作}

项目仍存在一些不足。首先，FER2013 灰度模型在 RAF-DB 外部测试上的泛化能力有限，说明低分辨率灰度数据与真实场景彩色数据之间存在明显域差异。其次，RAF-DB Compound 复合表情识别结果较低，主要受样本数量少、类别不平衡、类别语义接近和复合表情边界模糊影响。第三，RAF-ML 多标签任务虽然 Macro F1 较稳定，但固定阈值 0.5 未必适合所有情绪标签。第四，GUI 摄像头实时识别会受到光照、摄像头角度、人脸定位质量和用户距离影响。第五，当前系统主要基于表情类别标签进行监督学习，尚未进一步结合面部关键点、面部动作单元等可解释性信息。

后续工作可以从以下方向继续改进：

\begin{enumerate}
  \item \textbf{Compound 任务改进：}可以尝试类别重采样、focal loss、类别权重进一步调节和更强数据增强，以缓解小样本过拟合；也可以设计层次分类策略，先识别基础情绪成分，再细分复合类别，从而降低 11 类固定复合标签之间的直接竞争难度。
  \item \textbf{RAF-ML 多标签改进：}可以为不同情绪标签设置动态阈值或类别自适应阈值，而不是统一使用 0.5；还可以建模标签相关性，分析 Happiness--Surprise、Fear--Sadness 等情绪成分的共现关系，从而提高多标签识别的 Recall 和 Macro F1。
  \item \textbf{GUI 摄像头改进：}可以加入时间平滑或滑动窗口投票，避免连续帧预测结果频繁跳动；同时可以优化人脸检测失败时的中心区域回退策略，并在低置信度场景下给出更稳定的提示。
  \item \textbf{模型结构改进：}可以尝试 ResNet18、MobileNet、EfficientNet 或 Vision Transformer 等更强 backbone，但需要在准确率、参数量、推理速度和摄像头实时性之间进行权衡。
  \item \textbf{可解释性改进：}可以考虑引入面部关键点、面部动作单元或注意力热力图等信息，分析模型关注的面部区域是否符合人类对表情变化的理解，从而增强系统解释性和鲁棒性。
  \item \textbf{视频级建模：}可以进一步使用视频序列建模方法，利用表情随时间变化的动态信息，减少单帧识别中光照、模糊和瞬时姿态变化带来的不稳定。
\end{enumerate}

\section{总结}

本项目围绕“基于 PyQt5 与深度学习的人脸表情识别系统设计与实现”这一主题，完成了从数据集整理、模型设计、训练评估、消融对比到桌面端 GUI 展示的完整流程。项目并非只实现单一分类模型，而是以人脸表情识别任务为主线，逐步覆盖 FER2013 灰度基础实验、CK+ 和 RAF-DB 外部泛化测试、RAF-DB RGB 真实场景模型训练、AffectNet RGB Full 训练、AffectNet 预训练迁移、RAF-DB Compound 复合表情识别以及 RAF-ML 多标签混合表情识别等多个实验阶段。通过这些阶段，报告较系统地比较了输入形式、数据来源、训练策略和任务定义对模型性能的影响。

从模型结果看，灰度 FER2013 模型能够作为基础表情识别基线，但在 RAF-DB 等真实场景彩色数据上的外部泛化能力有限；RAF-DB RGBMiniResNet 在同源测试集上取得了较好的 Accuracy 和 Macro F1，说明 RGB 图像中的颜色、纹理和真实场景信息对表情识别具有明显帮助。AffectNet 预训练 + RAF-DB 微调模型取得了接近 RAF-DB 从头训练的性能，体现出预训练迁移的价值，但由于 AffectNet 与 RAF-DB 在采集来源、图像风格、标签映射和类别分布上存在差异，最终结果并未完全超过同源数据训练模型。Compound 复合表情识别的表现相对较低，反映出复合表情类别语义接近、样本数量较少和类别边界模糊等问题；相比之下，RAF-ML 多标签混合表情任务通过 sigmoid 输出多个情绪成分，更适合描述真实场景中复杂表情可能同时包含多种情绪的情况。

从系统实现看，项目将多个实验模型统一接入 PyQt5 GUI，支持图片上传识别和摄像头实时识别，并实现了基本表情识别、复合表情识别和混合表情识别三类任务的动态切换。GUI 能够根据任务类型自动调整标签列表、输出方式和概率条数量，避免将 RAF-ML 多标签模型误用为普通 softmax 分类模型。摄像头部分集成了 YuNet、MediaPipe 和中心区域识别三种人脸定位方式，使训练得到的模型能够以更直观的形式进行实时演示。消融分析脚本进一步将已有实验结果整理为表格和柱状图，使模型性能差异更易于解释。

总体而言，本项目较完整地体现了人工智能课程项目中“算法模型 + 实验分析 + 交互系统”的综合实践思路。模型层面完成了从基础 CNN 到轻量残差网络、从单标签分类到多标签识别的扩展；实验层面通过多数据集、多任务和消融对比分析了模型性能变化原因；系统层面通过 GUI 将离线训练结果转化为可演示应用。当然，项目仍存在进一步提升空间，例如可以使用更强的视觉 backbone、进一步缓解类别不平衡、引入面部关键点或动作单元增强解释性，并在视频序列层面建模表情的时间变化。这些方向可以作为后续改进和扩展研究的基础。

\section{参考资料与数据来源}

\begingroup
\small
\sloppy
\begin{enumerate}
  \item FER2013 ImageFolder Kaggle 数据集：\url{https://www.kaggle.com/datasets/astraszab/facial-expression-dataset-image-folders-fer2013}
  \item CK+ Kaggle ckdataset：\url{https://www.kaggle.com/datasets/davilsena/ckdataset}
  \item RAF-DB Kaggle 数据集：\url{https://www.kaggle.com/datasets/shuvoalok/raf-db-dataset}
  \item RAF 官方 Basic / Compound 页面：\url{http://www.whdeng.cn/RAF/model1.html}
  \item AffectNet YOLO Format Kaggle 数据集：\url{https://www.kaggle.com/datasets/fatihkgg/affectnet-yolo-format}
  \item RAF-ML 官方网站：\url{http://www.whdeng.cn/RAF/model2.html}
  \item OpenCV YuNet 模型来源：\url{https://github.com/opencv/opencv_zoo/tree/main/models/face_detection_yunet}
  \item MediaPipe Face Detector 文档：\url{https://ai.google.dev/edge/mediapipe/solutions/vision/face_detector}
  \item PyTorch 官方文档：\url{https://pytorch.org/docs/stable/index.html}
  \item PyQt5 文档：\url{https://www.riverbankcomputing.com/static/Docs/PyQt5/}
  \item OpenCV 官方文档：\url{https://docs.opencv.org/}
\end{enumerate}
\endgroup

\end{document}
